% summary.tex -- Module 8: Last-minute summary (the night before the exam).
% Ideas, formulas and traps are taken from the big ideas, formula cards, Common mistakes
% the other values are quoted from the chapters.
\unnumberedchapter[\faClipboardList]{Last-minute summary}

Read this the night before the exam. Formulas marked
\textcolor{hscmistake}{\bfseries learn} are \emph{not} on the formula sheet.

% a little extra room between lines that hold fractions
\setlength\lineskiplimit{1.5pt}\setlength\lineskip{2.5pt}
\newcommand\sumLearn{\textcolor{hscmistake}{\bfseries learn}}
\newcommand\sumhead[1]{\par\smallskip\noindent{\bfseries\color{hscbrand}#1}\enspace}
\newenvironment{sumlist}{\begin{itemize}[itemsep=0pt, topsep=2pt]}{\end{itemize}}

\begin{summarybox}[fontupper=\small]{\Cref{ch1:chap}: Origins of the Elements}
\begin{sumlist}
  \item Big Bang, 13.8 billion years ago: photons made pairs ($E \ge 2mc^2$). Quarks $\to$
        protons and neutrons $\to$ helium nuclei $\to$ atoms and the CMB.
  \item Evidence: galaxy redshift grows with distance (Hubble), the CMB (\qty{2.7}{K}),
        \qty{25}{\percent} helium.
  \item Hot dense object: continuous spectrum. Hot gas: bright lines. Cool gas in front:
        dark lines, at the same wavelengths, so lines identify the elements. Classes
        O B A F G K M, from hot (blue) to cool (red).
  \item HR diagram: hot stars on the \emph{left}, bright at the top. Mass decides a
        star's life: below $8\,M_\odot$ $\to$ white dwarf; above $\to$ supernova.
  \item p--p chain and CNO cycle: 4 protons $\to$ helium-4 $+$ \qty{26.7}{MeV}. Giants
        and supergiants make heavier elements, up to iron.
\end{sumlist}
\sumhead{Formulas} On the sheet: $E = mc^2$, \ $\lambda_\mathrm{max} = \frac{b}{T}$.
\ \sumLearn: $z \approx \frac{v}{c}$, \ $v = H_0 d$, \ $t \approx \frac{1}{H_0}$, \
$\frac{L_A}{L_B} = 100^{\frac{M_B - M_A}{5}}$, \ $E\,(\mathrm{MeV}) = \Delta m\,(\mathrm{u})
\times 931.5$.
\sumhead{Traps} $T$ in kelvin. $H_0$ in \si{s^{-1}} before $\frac{1}{H_0}$. Smaller
magnitude $=$ brighter.
\end{summarybox}

\begin{summarybox}[fontupper=\small]{\Cref{ch2:chap}: Structure of the Atom}
\begin{sumlist}
  \item Electric and magnetic fields bend cathode rays (towards the $+$ plate), so they
        are negative particles, not waves.
  \item Thomson (1897): crossed fields give $v = \frac{E}{B}$; $B$ alone gives
        $\frac{q}{m}$, about 1800 times the hydrogen ion's. Electrons are in all atoms.
  \item Millikan (1909--1913): every oil-drop charge is a whole number times
        \qty{1.602e-19}{C}. Charge is quantised.
  \item Geiger--Marsden (1909): most alphas pass through gold foil, about 1 in 8000
        bounce back. So: a tiny, dense, positive nucleus (Rutherford, 1911).
  \item Chadwick (1932): momentum and energy conservation show a neutral particle with a
        proton's mass: the neutron.
\end{sumlist}
\sumhead{Formulas} On the sheet: $E = \frac{V}{d}$, \ $\vec{F} = q\vec{E}$, \
$F = qvB\sin\theta$, \ $F_\mathrm{c} = \frac{mv^2}{r}$, Coulomb's law, momentum and $K$
conservation. \ \sumLearn: $v = \frac{E}{B}$, \ $\frac{q}{m} = \frac{E}{B^2 r}$, \
$\frac{q}{m} = \frac{2V}{B^2 r^2}$, \ $q = \frac{mgd}{V}$, \ $q = ne$.
\sumhead{Traps} $v = \frac{E}{B}$, not $\frac{B}{E}$. Reverse the right-hand rule for
electrons.
\end{summarybox}

\begin{summarybox}[fontupper=\small]{\Cref{ch3:chap}: The Quantum Atom}
\begin{sumlist}
  \item Rutherford's atom should collapse and give a continuous spectrum. It does not.
  \item Bohr: stationary states do not radiate; one jump gives one photon,
        $hf = E_\mathrm{i} - E_\mathrm{f}$; hydrogen $E_n = -\frac{\qty{13.6}{eV}}{n^2}$.
  \item Lyman ($n_\mathrm{f} = 1$, UV); Balmer ($n_\mathrm{f} = 2$: 656, 486, 434,
        \qty{410}{nm}); Paschen ($n_\mathrm{f} = 3$, IR).
  \item De Broglie: $\lambda = \frac{h}{mv}$, confirmed by electron diffraction (1927).
  \item Bohr works for hydrogen only. Schr\"odinger: $\psi^2$ $=$ probability; orbitals,
        not orbits.
\end{sumlist}
\sumhead{Formulas} On the sheet: $E = hf$, \ $\frac{1}{\lambda} =
R\left(\frac{1}{n_\mathrm{f}^2} - \frac{1}{n_\mathrm{i}^2}\right)$, \
$\lambda = \frac{h}{mv}$. \ \sumLearn: $E = \frac{hc}{\lambda}$, \ $E_n$, \
$2\pi r = n\lambda$, \ $mvr = \frac{nh}{2\pi}$.
\sumhead{Traps} Use the \emph{gap} between levels. eV to J, nm to m. Square the $n$
values.
\end{summarybox}

\begin{summarybox}[fontupper=\small]{\Cref{ch4:chap}: Properties of the Nucleus}
\begin{sumlist}
  \item Alpha: helium-4 nucleus ($A - 4$, $Z - 2$), stopped by paper. Beta-minus:
        $\mathrm{n} \to \mathrm{p} + \mathrm{e}^- + \overline{\nu}_\mathrm{e}$ ($Z + 1$),
        stopped by aluminium. Gamma: a photon, only reduced by thick lead.
  \item Decay is random, but the amount halves every half-life, for any sample size.
  \item A nucleus has less mass than its parts: $E_\mathrm{B} = \Delta m\,c^2$. Binding
        energy per nucleon peaks at iron-56. Fission and fusion move towards iron, so
        mass is lost and energy is released.
  \item Fission: U-235 $+$ slow neutron $\to$ 2 fragments $+$ 2--3 neutrons $+$
        about \qty{200}{MeV}. Moderator slows neutrons; control rods absorb them.
        Fissions each second $= \frac{P}{E_\mathrm{fission}}$ ($P$ in W, $E_\mathrm{fission}$
        in J).
  \item Fusion needs a very hot, dense plasma to beat the electric repulsion.
\end{sumlist}
\sumhead{Formulas} On the sheet: $N_\mathrm{t} = N_0 e^{-\lambda t}$, \
$\lambda = \frac{\ln 2}{t_{\frac{1}{2}}}$, \ $E = mc^2$. \ \sumLearn: $\sum A$ and
$\sum Z$ conserved, \ $A = \lambda N$, \ $N_\mathrm{t} = N_0\left(\frac{1}{2}\right)^n$, \
$\Delta m = (Zm_\mathrm{p} + Nm_\mathrm{n}) - m_\mathrm{nucleus}$, \
$E\,(\mathrm{MeV}) = \Delta m\,(\mathrm{u}) \times 931.5$.
\sumhead{Traps} $\lambda$ and $t$ in the same unit. ``75\,\% decayed'' means 25\,\% is
left.
\end{summarybox}

\begin{summarybox}[fontupper=\small]{\Cref{ch5:chap}: Deep Inside the Atom}
\begin{sumlist}
  \item SLAC (1968): electrons bounced at large angles, so protons contain point-like
        quarks. New particles: positron, muon, pions, neutrino.
  \item 6 quarks ($+\frac23 e$ or $-\frac13 e$), 6 leptons ($-e$ or 0); bosons: gluon,
        photon, W, Z, Higgs. Baryons: 3 quarks (p is uud). Mesons: quark $+$ antiquark.
  \item Electric fields speed particles up; magnets bend them. Linac, cyclotron,
        synchrotron (LHC). Found: W and Z (1983), top (1995), Higgs (2012).
  \item Still incomplete: gravity, dark matter, dark energy, neutrino mass, antimatter.
\end{sumlist}
\sumhead{Formulas} On the sheet: $\lambda = \frac{h}{mv}$, \ $p_\mathrm{v}$, \
$W = qV$, \ $F = qvB\sin\theta$. \ \sumLearn: $Q = q_1 + q_2 + q_3$, \
$r = \frac{mv}{qB}$, \ $T = \frac{2\pi m}{qB}$.
\sumhead{Traps} Flip the sign for an antiquark. Mesons are bosons. Magnets do no work.
\end{summarybox}
